problem:  
Let \((M^n,g,f)\) be a complete, noncompact **gradient shrinking Ricci soliton** satisfying  
\[
\mathrm{Ric}_g + \nabla^2 f = \tfrac{1}{2} g, \qquad \min f = 0,
\]
with **nonnegative sectional curvature** and **Euclidean volume growth**.  
Endow \(M\) with the normalized weighted measure  
\[
d\mu_f := (4\pi)^{-n/2} e^{-f}\, dV_g,
\]
and for any smooth domain \(\Omega \subset M\), define  
\[
\mu_f(\Omega) = \int_\Omega d\mu_f, \quad 
\mu_f(\partial \Omega) = \int_{\partial \Omega} e^{-f} \, dA_g.
\]
The **weighted isoperimetric profile** is the function
\[
I_f(V) = \inf_{\Omega:\, \mu_f(\Omega) = V}\, \mu_f(\partial \Omega), \quad V \in (0,1),
\]
defined with respect to the normalized measure \(\mu_f\).

Every shrinking soliton satisfies the Bakry–Émery curvature lower bound \(\mathrm{Ric}_f = \mathrm{Ric}_g + \nabla^2 f = \tfrac{1}{2} g\), i.e. the synthetic curvature–dimension condition \(CD(\tfrac{1}{2},\infty)\).  
In Gaussian space \((\mathbb{R}^n,g_{\mathrm{Euc}},f=|x|^2/4)\), the function \(I_f(V)\) is **strictly concave**, and all isoperimetric minimizers are Euclidean halfspaces.

**Problem:**  

1. (**Rigidity via strict concavity**)  
For dimensions \(n\le4\), determine whether the **strict concavity of \(I_f(V)\)** uniquely characterizes the Gaussian shrinker among all noncompact gradient shrinking solitons with nonnegative sectional curvature and Euclidean volume growth.  
That is,  
> If \((M,g,f)\) satisfies \(CD(\tfrac{1}{2},\infty)\) and \(I_f''(V)<0\) for all regular \(V \in (0,1)\), must \((M,g,f)\) be isometric (up to scaling) to the Gaussian shrinker?

2. (**Quantitative concavity‑defect**)  
Define the **concavity‑defect functional**
\[
\tilde{\mathcal D}_f := 
\sup_{0<V_1<V<V_2<1}
\left[
I_f(V) - 
\frac{(V_2-V)\, I_f(V_1) + (V-V_1)\, I_f(V_2)}{V_2 - V_1}
\right].
\]
This functional measures the maximal deviation of \(I_f\) from concavity.  

> Does \(\tilde{\mathcal D}_f = 0\) hold if and only if \((M,g,f)\) is the Gaussian shrinker?  
> Moreover, can the size or asymptotic behaviour of \(\tilde{\mathcal D}_f\) (e.g. for small \(V\)) be expressed geometrically—for instance, via curvature integrals, eigenvalue data of the drift Laplacian, or the structure of the soliton's asymptotic cone?

---

Why is it a "good" problem:  
This question links the **geometry of Ricci shrinkers** with the **isoperimetric theory of metric‑measure spaces** under curvature‑dimension conditions. It refines the relationship between self‑similar Ricci flow geometry and measure‑concentration inequalities into a concrete analytic test: *does concavity of the weighted isoperimetric profile rigidly identify the Gaussian shrinker?*

- **Novel perspective:**  
  While \(CD(K,\infty)\) spaces enjoy generalized convexity and concentration inequalities, the *strict* concavity of \(I_f(V)\) is not guaranteed by these synthetic conditions. It could act as a new rigidity criterion distinguishing the Gaussian soliton from all others.

- **Bridge between fields:**  
  It draws connections among **Ricci soliton geometry**, **Bakry–Émery analysis**, and **optimal transport–based isoperimetric inequalities**, requiring cross‑field methods from geometric PDE, spectral theory, and geometric measure theory.

- **Simplicity with depth:**  
  The entire problem is built from tangible quantities—volume, perimeter, concavity—yet its resolution involves subtle analysis of the soliton potential, curvature decay, and functional inequalities, echoing Perelman‑type monotonicity phenomena in a static setting.

- **Low‑dimensional focus:**  
  The restriction \(n\le4\) reflects the frontier of current classification results while leaving the isoperimetric characterization of these shrinkers completely open.

- **Potential impact:**  
  A *positive rigidity result* would yield a purely analytic characterization of the Gaussian shrinker via a measure‑theoretic property intrinsic to its geometry, forging a new link between Ricci‑flow models and probabilistic isoperimetry.  
  Conversely, a *negative result* would uncover novel non‑Gaussian examples where the expected concavity fails, shedding light on the limitations of curvature‑dimension methods in capturing soliton rigidity.

This combination of simplicity, depth, and cross‑disciplinary synthesis makes it a genuinely “good” research problem.